\begin{abstract}
We give a diagrammatic representation of deductive and inductive reasoning: deduction is a certain
passage from thought to existence, $\varphi\to\exists$, and induction a tentative, iterated passage
from the study of existence to the theory of thought, $s\exists\circlearrowright t\varphi$. We make
the diagram precise in a possible-worlds semantics, in which deduction is semantic entailment and a
theory is supported by a study when some world is consistent with both. We prove that deduction is
sound and transitive, that iterated observation can only narrow the supported theories, that a
single counterexample falsifies a universal theory, and that no finite study entails a universal
theory over an unobserved individual, which is the logical part of Hume's problem of induction. A
Bayesian refinement shows that every genuine test a theory predicts confirms it; a ranked semantics
treats defeasible inference; and ordered levels show iteration moving upward. We apply the
framework to plant disease diagnosis and forecasting: diagnosis is elimination and refinement,
Koch's postulates are hypothetico-deductive confirmation, forecasting rules are defeasible, and a
new theorem characterizes when an assay can resolve a diagnosis. A quantum information model
realizes the same structure, with the classical setting as the case of commuting observables and
the single-basis limit of state tomography as an instance of assay resolution. Every theorem is
verified in the Lean~4 proof assistant with Mathlib.
\end{abstract}
